\section{Cosets, flags and residue generation}\label{O-sec:cosets}
All type sets in this paper are finite. Cosets are left cosets. In a coset geometry $(B,(B_i)_{i\in I})$ write $B_J=\bigcap_{i\in J}B_i$, with $B_\varnothing=B$. A regular hypertope is thin, residually connected and flag-transitive. Automorphisms preserve types. In its faithful standard model, the group acts freely and transitively on chambers. Its distinguished involutions $b_i$ interchange a fixed chamber and its $i$-adjacent chamber; then
\begin{equation}\label{O-eq:parabolic}
B_J=\langle b_i:i\in I\setminus J\rangle,\qquad B_I=1.
\end{equation}
We recall why the geometric hypotheses give these assertions. In a thin geometry a type-preserving automorphism fixing a chamber fixes all its adjacent chambers. Residual connectivity implies that the chamber graph, and the chamber graph of every residue, is connected: by induction on rank, chambers containing a given element are connected; a path in the incidence graph connects these families, using chambers containing each consecutive incident pair. The rank-one and rank-two cases start the induction. Thus a chamber stabilizer is trivial. The same argument inside the residue of the base flag of types $J$ proves (\ref{O-eq:parabolic}). It also proves that the full type-preserving automorphism group is the displayed group whenever a faithful group already acts transitively on chambers.

For a second regular hypertope $(A,(A_x)_{x\in S})$, put $A(T)=\langle a_x:x\in T\rangle$ for $T\subseteq S$. Formula (\ref{O-eq:parabolic}) gives
\begin{equation}\label{O-eq:ip}
A(T)\cap A(T')=A(T\cap T').
\end{equation}
We will use both this intersection property and flag-transitivity. They are different requirements: the intersection property alone is not a general criterion for a nonstring coset system to be a geometry.

\begin{lemma}\label{F-lem:cosets}
A coset incidence system $(A,(A_i)_{i\in I})$ is flag-transitive if
and only if the family of all left cosets of all $A_J$, $J\subseteq I$,
is Helly.
\end{lemma}
\begin{proof}
A flag is a collection of pairwise intersecting cosets of distinct
maximal parabolics. It is a translate of a standard flag precisely
when these cosets have a common element. This proves one implication.
For the other, expand each coset $aA_J$ into the intersection of
$aA_j$, $j\in J$. If two original cosets intersect, every maximal
coset from one meets every maximal coset from the other. Within one
expanded coset they meet at $a$. Two such cosets of the same maximal
parabolic are therefore equal. The resulting distinct-type collection
is a flag, and flag transitivity supplies a common element. Empty
$J$ imposes no constraint. A standard flag extends to the standard
chamber, so flag transitivity also supplies the geometric property.
\end{proof}

\begin{lemma}[Coset intersection]\label{O-lem:helly}
If $B$ is flag-transitive on its coset geometry, every finite pairwise-intersecting family of cosets of standard parabolic subgroups has nonempty total intersection. In particular, for $H,K,J\subseteq I$,
\[
B_HB_J\cap B_KB_J=B_{H\cup K}B_J.
\]
\end{lemma}
\begin{proof}
Lemma~\ref{F-lem:cosets} gives the first assertion. For the last assertion, if $g\in B_HB_J\cap B_KB_J$, the three cosets $B_H,B_K,gB_J$ intersect pairwise. A point of their common intersection belongs to $B_{H\cup K}$ and to $gB_J$. The reverse inclusion follows from subgroup containment.
\end{proof}

\begin{theorem}[Element orbits]\label{O-thm:elements}
Fix $j\in I$ and $x_0=B_j\in X=B/B_j$. If $B$ is flag-transitive, then
\begin{equation}\label{O-eq:orbit}
(B_Hx_0)\cap(B_Kx_0)=B_{H\cup K}x_0\qquad(H,K\subseteq I).
\end{equation}
Thus the complementary notation $O(M)=B_{I\setminus M}x_0$ satisfies $O(M)\cap O(N)=O(M\cap N)$.
\end{theorem}
\begin{proof}
The inverse image of $B_Hx_0$ under $B\to B/B_j$ is $B_HB_j$. Apply Lemma~\ref{O-lem:helly} and take images. Finally $(I\setminus M)\cup(I\setminus N)=I\setminus(M\cap N)$.
\end{proof}





\subsection{Parabolics with a nontrivial chamber stabilizer}
We now allow flag-transitive coset geometries whose acting groups are not
faithful. The type sets remain finite and nonempty. In a geometry
$\alpha=(A,(A_x)_{x\in S})$ put $A_E=\bigcap_{x\in E}A_x$, with
$A_\varnothing=A$, and call $A_S$ the Borel subgroup. It need not be
trivial. The residue of the standard flag of types $E$ is the coset
geometry of $A_E$ with parabolics $A_{E\cup\{x\}}$, $x\notin E$.
To see this identification, take an element incident with that flag;
flag transitivity supplies a common representative in $A_E$. Two such
representatives define the same element exactly when their quotient
lies in $A_E\cap A_x$. This also identifies incidence by intersecting
the corresponding cosets.

For a coset geometry of rank at least two the incidence graph is
connected if and only if its maximal parabolics generate the acting
group. Indeed, if $H$ is that generated subgroup, paths from the
standard flag remain among cosets with representatives in $H$.
Conversely multiplication successively by elements of the parabolics
provides paths through incident cosets; when two successive steps have
the same type, an intervening different type gives the same conclusion.
Thus the components are the translates of this component. Applied to
every residue, this proves that residual connectedness is equivalent to
\begin{equation}\label{O-eq:rc-borel}
 A_E=\langle A_{E\cup\{x\}}:x\notin E\rangle
 \quad\text{if }|S\setminus E|\geq2.
\end{equation}
For rank one residual connectedness imposes no condition.

\begin{lemma}[Generation of parabolics]\label{O-lem:parabolic-join}
If $\alpha$ is residually connected, then
\[
 A_E=\langle A_S,\ A_{S\setminus\{x\}}:x\notin E\rangle.
\]
Consequently, for any nonempty family $(E_\nu)$ of subsets of $S$,
\begin{equation}\label{O-eq:parabolic-join}
 \langle A_{E_\nu}:\nu\rangle=A_{\bigcap_\nu E_\nu}.
\end{equation}
\end{lemma}
\begin{proof}
Induct on $|S\setminus E|$. For zero this is $A_S=A_S$; for one
the displayed generator is $A_E$ itself. For larger complements,
use (\ref{O-eq:rc-borel}) and the induction hypothesis on each child.
Every missing type remains missing in at least one child, so all and
only the indicated minimal parabolics occur. For the second assertion,
each group on the left lies in the right-hand group. Conversely, every
minimal parabolic required to generate the right-hand group occurs
in some $A_{E_\nu}$, and $A_S$ lies in every one. The nonempty-family
hypothesis is necessary when $A_S\ne1$.
\end{proof}
The explicit $A_S$ in the first formula must not be dropped at
$E=S$. This avoids the empty-generated-subgroup problem in the
source's background equivalences.
